%%% Преамбула авторского отчета %%%
\usepackage[T2A]{fontenc}
\usepackage[utf8]{inputenc}
\usepackage[russian]{babel}
\usepackage{graphicx}
\usepackage{wrapfig}
\usepackage{color,xcolor}
\usepackage{booktabs}
\usepackage{multirow}
\usepackage{tabularx}
\usepackage{caption}
\usepackage{subcaption}

% Шрифты и оформление текста
\usepackage{tempora}
\usepackage{mathtools,amssymb,amsmath}
\usepackage{newtxmath}
\usepackage{indentfirst}
\usepackage{float}
\usepackage{enumitem}

% Геометрия страницы
\usepackage{geometry}
\geometry{left=2.5cm, right=1.5cm, top=2cm, bottom=2cm}

\renewcommand{\baselinestretch}{1.3}
\parindent 1.25cm
\sloppy
\hyphenpenalty=1000
\clubpenalty=10000
\widowpenalty=10000

% Настройка списков
\setlist[enumerate,itemize]{leftmargin=1.25cm, topsep=2pt, itemsep=2pt, parsep=0pt}
\renewcommand{\labelitemi}{--}

% Гиперссылки
\usepackage{hyperref}
\hypersetup{
    colorlinks=true,
    linkcolor=blue!70!black,
    citecolor=blue!70!black,
    urlcolor=blue!70!black,
    pdfauthor={Савватей Степанов},
    pdftitle={Исследование методов планирования последовательностей интенций поверх Forward-Backward представлений}
}

% TikZ графика
\usepackage{tikz}
\usetikzlibrary{shapes,arrows,positioning,calc,fit,backgrounds,decorations.pathreplacing,patterns}

% Блоки алгоритмов
\usepackage[linesnumbered,ruled,vlined]{algorithm2e}
\renewcommand{\algorithmcfname}{Алгоритм}
\SetKwInput{KwIn}{Вход}
\SetKwInput{KwOut}{Выход}
\SetKwInput{KwData}{Данные}
\SetKwInput{KwResult}{Результат}
\SetKw{KwAnd}{и}
\SetKw{KwOr}{или}

% Команды для формул и секций (в стиле диплома)
\newcommand{\addequation}[2]{
    \begin{equation}
        \label{#2}
        #1
    \end{equation}
}

\newcommand{\anonsection}[1]{
    \section*{#1}
    \addcontentsline{toc}{section}{#1}
}

\numberwithin{equation}{section}
\counterwithin{figure}{section}
\counterwithin{table}{section}
